\chapter*{Agradecimientos}

Deseo expresar mi más profundo agradecimiento a mi profesor guía, el Dr.~Alejandro Weinstein, por su dirección rigurosa, su paciencia inagotable y su capacidad para transformar cada obstáculo en una oportunidad de aprendizaje. Su visión sobre la intersección entre procesamiento de señales y neurociencia ha sido una fuente constante de inspiración a lo largo de todo este proceso.

Al Dr.~Matías Zañartu, profesor co-guía de esta tesis, le agradezco su mirada crítica y su insistencia en que cada afirmación debía estar respaldada por evidencia. Sus preguntas —a veces incómodas, siempre necesarias— mejoraron sustancialmente la calidad de este trabajo.

Al Departamento de Ingeniería Electrónica de la Universidad Técnica Federico Santa María, por proporcionar el entorno académico y los recursos computacionales que hicieron posible esta investigación.

A mis compañeros del laboratorio, por las discusiones técnicas a altas horas de la noche y por recordarme que la investigación también es comunidad. A los desarrolladores de MNE-Python, Optuna, PyGSP y el ecosistema de código abierto científico de Python, cuyo trabajo colectivo hizo posible que esta tesis se construyera sobre hombros de gigantes.

Finalmente, a mi familia —a mis padres, por enseñarme que la curiosidad es el motor más noble; a mi pareja, por su apoyo incondicional durante los meses de escritura intensiva; y a mis amigos, por entender mis ausencias y celebrar mis pequeños avances como si fueran propios.

\vspace{2em}
\begin{flushright}
\textit{\ThesisCity, \ThesisMonth\ de \ThesisYear}
\end{flushright}